% Part: first-order-logic
% Chapter: introduction
% Section: formulas

\documentclass[../../../include/open-logic-section]{subfiles}

\begin{document}

\olfileid[fr]{fol}{int}{fml}

\section{\usetoken{P}{formula}}

Voici la méthode que nous suivrons pour définir rigoureusement les
!!{sentence}s de la logique du premier ordre et traiter les difficultés
liées à l'emploi des !!{variable}s. Nous commençons par définir un
ensemble \emph{différent} d'expressions : les !!{formula}s. Nous
pourrons ensuite examiner le rôle qu'y jouent les !!{variable}s et, à
partir des notions de !!{variable} « libre » et « liée », définir ce
qu'est !!a{sentence}, à savoir !!a{formula} sans !!{variable} libre.
Cette démarche ne rend pas seulement la définition d'« !!{sentence} »
plus maniable : elle sera aussi essentielle à notre définition de la
notion sémantique de satisfaction.

Définissons « !!{formula} » pour un langage du premier ordre simple,
qui ne contient qu'!!a{predicate}~$\Obj P$, !!a{constant}~$\Obj a$ et
les seuls symboles logiques $\lnot$, $\land$ et~$\lexists$. Nos
définitions complètes seront beaucoup plus générales : elles
autoriseront une infinité de !!{predicate}s et de !!{constant}s. Nous
considérerons aussi des !!{function}s qui, combinés avec des
!!{constant}s et des !!{variable}s, forment des « termes ». Pour
l'instant, $\Obj a$ et les variables seront nos seuls termes. Il nous
faut une infinité de !!{variable}s. Nous prendrons officiellement pour
variables les symboles $\Obj v_0$, $\Obj v_1$, \dots.

\begin{defn}
L'ensemble des \emph{!!{formula}s}~$\Frm$ est défini comme suit :
\begin{enumerate}
\item\ollabel{fmls-atom} $\Atom{\Obj P}{\Obj a}$ et
  $\Atom{\Obj P}{\Obj v_i}$ sont des !!{formula}s ($i \in \Nat$).

\tagitem{prvNot}{\ollabel{fmls-not}Si $!A$ est !!a{formula}, alors
  $\lnot !A$ est !!a{formula}.}{}

\tagitem{prvAnd}{Si $!A$ et $!B$ sont des !!{formula}s, alors
  $(!A \land !B)$ est !!a{formula}.}{}

\tagitem{prvEx}{\ollabel{fmls-ex}Si $!A$ est !!a{formula} et si $x$
  est !!a{variable}, alors $\lexists[x][!A]$ est !!a{formula}.}{}

\tagitem{limitClause}{\ollabel{fmls-limit}Rien d'autre n'est
  !!a{formula}.}{}
\end{enumerate}
\end{defn}

La clause~\olref{fmls-atom} nous dit que $\Atom{\Obj P}{\Obj a}$ et
$\Atom{\Obj P}{\Obj v_i}$ sont des !!{formula}s pour tout
$i \in \Nat$. Ce sont les !!{formula}s dites \emph{atomiques}; elles
fournissent notre point de départ. Les autres clauses donnent des
moyens de former de nouvelles !!{formula}s à partir de celles qui sont
déjà formées. Ainsi, la clause~\olref{fmls-not} montre que
$\lnot \Atom{\Obj P}{\Obj v_2}$ est !!a{formula}, puisque
$\Atom{\Obj P}{\Obj v_2}$ est déjà !!a{formula} d'après la
clause~\olref{fmls-atom}. La clause~\olref{fmls-ex} montre ensuite que
$\lexists[\Obj v_2][\lnot \Atom{\Obj P}{\Obj v_2}]$ est une autre
!!{formula}, et ainsi de suite. La clause~\olref{fmls-limit} affirme que
\emph{seules} les chaînes obtenues de cette manière sont des
!!{formula}s. En particulier,
$\lexists[\Obj v_0][\Atom{\Obj P}{\Obj a}]$ et
$\lexists[\Obj v_0][\lexists[\Obj v_0][\Atom{\Obj P}{\Obj a}]]$
\emph{sont} des !!{formula}s, tandis que
$(\lnot \Atom{\Obj P}{\Obj a})$ n'en est pas une, à cause des
parenthèses extérieures superflues.

Cette manière de définir les !!{formula}s s'appelle une
\emph{définition inductive}; elle permet de démontrer des propriétés des
!!{formula}s au moyen d'une variante du raisonnement par induction
appelée \emph{induction structurelle}. Ces deux notions sont présentées
de façon générale dans les \olref[mth][ind][idf]{sec} et
\olref[mth][ind][sti]{sec}, qu'il convient de revoir avant d'aborder
les démonstrations ultérieures. L'idée est la suivante. Pour montrer
qu'une propriété vaut pour toutes les !!{formula}s, on commence par
montrer qu'elle vaut pour les !!{formula}s atomiques. On montre ensuite
que, \emph{si} elle vaut pour une !!{formula}~$!A$ (et pour~$!B$),
elle vaut \emph{aussi} pour $\lnot !A$, $(!A \land !B)$ et
$\lexists[x][!A]$. Cela établit par exemple la propriété pour
$\lexists[\Obj v_2][\lnot \Atom{\Obj P}{\Obj v_2}]$ : la première
étape la donne pour la !!{formula} atomique
$\Atom{\Obj P}{\Obj v_2}$; la deuxième la transmet à
$\lnot \Atom{\Obj P}{\Obj v_2}$, puis à
$\lexists[\Obj v_2][\lnot \Atom{\Obj P}{\Obj v_2}]$ elle-même.
Puisque toutes les !!{formula}s sont engendrées inductivement à partir
des !!{formula}s atomiques, le raisonnement vaut pour chacune d'elles.

\end{document}
